% !PRINTED 312
\weberpage{312}
%<<CONT
from which it is to be concluded:
\weq{7}{\begin{aligned} v_\infty &= v_0 = \frac{-1+\sqrt5}{2}, \\ v_1 &= v_4 = \frac{-1-\sqrt5}{2}. \end{aligned}}
($v_0$, $v_\infty$ must be positive, since $q$ is real in them.)

From (3) there then results:
\weq{8}{\begin{aligned} w_0 &= 0, \quad w_1 = w_3 = -5-2i\sqrt5, \\ w_2 &= w_4 = -5+2i\sqrt5. \end{aligned}}
These must be the roots of (5) for $j=12^3$, and from this one obtains the resolvent sought in the form:
\weq{9}{w(w^2+10w+45)^2 = j(\omega) - 12^3 = \gamma_3^2.}
This equation can also be brought into the form:
\weq{10}{(w+3)^3(w^2+11w+64) = j(\omega) = \gamma_2^3,}
which one also reaches directly if, in place of setting $\omega=i$,
\[ \omega = \frac{-1+i\sqrt3}{2} \]
is assumed.

The resolvent found simplifies still further if [according to (9)]:
\[ z = \sqrt{w} = \frac{\gamma_3}{w^2+10w+45} \]
is set:
\weq{11}{z^5+10z^3+45z = \gamma_3,}
or if [according to (10)]
\weq{12}{w^2+11w+64 = y^3, \quad y = \frac{\gamma_2}{w+3}}
is set:
\weq{13}{y^5 - 40y^2 - 5\gamma_2 y - \gamma_2^2 = 0,}
an equation which also results if one makes the assumption from the outset:
\[ 2y = v_\infty v_0 + v_1 v_4 + v_2 v_3. \]

In equation (13), as one sees, the third and the fourth power of the unknown are missing. The same property belongs also, as is easily shown, to the equation whose roots
%CONT>>
